
\documentclass[12pt]{article}

\usepackage{amsmath, amsthm, amssymb, mathtools}
\usepackage{geometry}
\usepackage{xr-hyper}
\externaldocument{general-theory}[general-theory.pdf]
\externaldocument{convolution-equations}[convolution-equations.pdf]
\usepackage{hyperref}
\usepackage{enumitem}
\usepackage{microtype}
\usepackage{xcolor}

\geometry{margin=1.25in}

\newtheorem{theorem}{Theorem}[section]
\newtheorem{proposition}[theorem]{Proposition}
\newtheorem{lemma}[theorem]{Lemma}
\newtheorem{corollary}[theorem]{Corollary}
\newtheorem{conjecture}[theorem]{Conjecture}
\theoremstyle{definition}
\newtheorem{definition}[theorem]{Definition}
\newtheorem{example}[theorem]{Example}
\theoremstyle{remark}
\newtheorem{remark}[theorem]{Remark}
\newtheorem{warning}[theorem]{Warning}

\DeclareMathOperator{\supp}{supp}
\DeclareMathOperator{\type}{type}
\DeclareMathOperator{\Stab}{Stab}
\newcommand{\R}{\mathbb{R}}
\newcommand{\C}{\mathbb{C}}
\newcommand{\Z}{\mathbb{Z}}
\newcommand{\N}{\mathbb{N}}
\newcommand{\T}{\mathcal{T}}
\newcommand{\PW}{\mathrm{PW}}
\newcommand{\Sph}{\mathbb{S}}
\newcommand{\wh}{\widehat}
\newcommand{\abs}[1]{\left|#1\right|}
\newcommand{\norm}[1]{\left\|#1\right\|}
\newcommand{\ip}[2]{\left\langle #1,\, #2\right\rangle}
\newcommand{\re}{\operatorname{Re}}

% Status markers
\newcommand{\proven}{\textbf{\color{teal}[Proved]}}
\newcommand{\conditional}{\textbf{\color{orange}[Conditional]}}
\newcommand{\open}{\textbf{\color{red}[Open]}}
\newcommand{\heuristic}{\textbf{\color{purple}[Heuristic]}}
\newcommand{\review}[1]{\par\smallskip\noindent\textbf{\color{red}[Review comment.]} {\color{red}#1}\par\smallskip}
\newcommand{\response}[1]{\par\noindent\textbf{\color{blue}[Response.]} {\color{blue}#1}\par\smallskip}
\newcommand{\resolved}{\textbf{\color{teal}[RESOLVED]}}
\newcommand{\pending}{\textbf{\color{orange}[PENDING]}}
\newcommand{\fatal}{\textbf{\color{red}[OPEN]}}

\title{\textbf{Attempts and Exploratory Routes}\\[0.5em]
\large Failed, Partial, and Out-of-Reach Strategies}
\author{}
\date{\today}

\begin{document}

\maketitle
\tableofcontents
\bigskip

This file collects exploratory routes, higher-dimensional extensions, and
partial strategies that are currently out of reach or have not led to a
closed proof.  It is kept separate from the main $S^1$ sharp-constant
argument so that the latter remains focused.



% =======================================================================
\section{Abstract Convolution Side Routes}
\label{sec:abstract-side-routes}
% =======================================================================

We first record the parts of the abstract convolution discussion which are
useful as orientation, but do not themselves give the desired proof.  They
explain why the Fredholm route has to use the equation
$u=c(u^k\ast f)$ rather than the formally nearby equation
$u^k=c(u\ast f)$, and why a direct contraction argument is too local to
replace the spectral analysis.

\subsection{The equation $u^k=c\,(u\ast f)$: spectral expansion kills solutions}

In Fourier space the equation reads
$\wh u^{\ast k}(\xi)=c\,\wh u(\xi)\wh f(\xi)$.

\begin{proposition}[Type obstruction]  \label{prop:type-obstruction}
  Let $u\in\PW_\tau$ with $\tau>0$ satisfy $u^k=c(u\ast f)$ for $k\geq 2$.
  Then no positive-type solution is possible; every solution has type zero.
\end{proposition}

\begin{proof}
  The Titchmarsh convolution theorem gives
  $\supp(\wh u^{\ast k})\subseteq[-k\tau,k\tau]$, while the right side is
  supported in $[-\tau,\tau]$.  For $k\geq 2$ this forces $\tau=0$, so $u$
  is a polynomial.  Degree comparison then forces $u$ to be constant, and
  the remaining constants are determined by the scalar equation obtained
  from $u^k=c(u\ast f)$.
\end{proof}

\begin{corollary}
  For $k\geq 2$ the solution set of $u^k=c(u\ast f)$ in the Paley--Wiener
  class is a finite discrete set of constants.
\end{corollary}

This is a useful negative result, but it is not a route to the restriction
problem.  The Euler--Lagrange equation has the smoothing/compactness
orientation $u=c(u^k\ast f)$, not the type-expanding orientation above.

\subsection{The small-ball contraction argument}

For the compactly oriented equation
\[
  u=c(u^k\ast f),
\]
one can bypass Fredholm theory in a small ball.  If
$T(u):=c(u^k\ast f)$ and $u,v$ lie in a ball of radius $R$ in
$\PW_{\tau_f}$, Young's inequality gives
\[
  \norm{T(u)-T(v)}
  \leq \abs{c}kR^{k-1}\norm{f}_{L^1}\norm{u-v}.
\]
Thus $T$ is a contraction whenever
\[
  R<(\abs{c}k\norm{f}_{L^1})^{-1/(k-1)}.
\]
This gives local uniqueness near sufficiently small solutions, but it does
not control the nonlinear eigenfunctions relevant to the sharp restriction
constant.  In particular, it gives no quantitative isolation around the
constant solution at the physical normalisation, and it does not separate
the constant eigenvalue from possible nonconstant branches.

\subsection{Comparison of the two equations}

\renewcommand{\arraystretch}{1.5}
\begin{table}[h]
\centering
\begin{tabular}{lll}
\hline
& $u^k=c(u\ast f)$ & $u=c(u^k\ast f)$ \\
\hline
Fourier identity & $\wh u^{\ast k}=c\wh u\wh f$ &
  $\wh u=c\wh u^{\ast k}\cdot\wh f$ \\
Type constraint ($k\geq 2$) & Forces $\tau=0$ (Titchmarsh) &
  Allows $\tau\leq\tau_f$ \\
Operator structure & No compactness & $T$ compact \\
Linearisation & Not Fredholm in general &
  $I-K$, Fredholm index $0$ \\
Solutions ($k\geq 2$, generic) &
  Finitely many constants & Isolated; Leray--Schauder applies \\
\hline
\end{tabular}
\caption{Analytic contrast between the two convolution equations.}
\end{table}

The table is useful diagnostically: it shows that the equation with the
most rigid support behaviour is the wrong model for the restriction
Euler--Lagrange problem, while the correct model has Fredholm structure but
requires the harder compactness and spectral-gap inputs.



% =======================================================================
\section{The Higher-Dimensional Transversality Route}
\label{sec:higher-dim}
% =======================================================================

For $n\geq 2$ and the equation $\rho=cT(\rho)$ on $L^2(S^{n-1})$,
Proposition~\ref{prop:support-forcing} forces $\wh u$ onto $S^{n-1}$,
and the problem reduces to the incidence-fibre equation~\eqref{eq:sphere-eq}.
The symmetry group is $G=\R^n\rtimes SO(n)$, acting on densities by
$(x_0,R)\cdot\rho(\omega)=e^{-ix_0\cdot\omega}\rho(R^{-1}\omega)$.

\subsection{The five-step programme}

The route to finiteness of shapes requires five analytic inputs.  We state
each and assess its status for general~$n$.

\paragraph{Step 1: Smoothing / Fredholm estimate.}  \open\quad
The target is a multilinear smoothing estimate
\[
  A_k:H^s(S^{n-1})^k\to H^{s+\varepsilon}(S^{n-1}),
  \qquad
  \norm{A_k(\rho_1,\ldots,\rho_k)}_{H^{s+\varepsilon}}
  \leq C_s\prod_{j=1}^k\norm{\rho_j}_{H^s},
\]
for some $\varepsilon>0$, where $A_k$ is the $k$-linear operator behind
the incidence-fibre equation.  A weaker but sufficient target is
compactness of $A_k$ on bounded sets.

For the $S^1$ problem, the truncation argument
(Theorem~\ref{thm:K-compact}) bypasses this step entirely.  For $n\geq 2$,
the proof must exploit the curvature of $S^{n-1}$ and the transversality
of the constraint $\omega_1+\cdots+\omega_k=\omega$.  The key obstruction
is the existence of \emph{degenerate fibres}; see
Section~\ref{sec:degenerate-fibres}.

\paragraph{Step 2: Regularity bootstrap.}  \open\quad
Show that weak $L^2$ solutions are automatically smooth:
$\rho\in L^2(S^{n-1}),\;\rho=cT(\rho)\implies\rho\in H^s$ for all~$s$.
If Step~1 gives $T:L^2\to H^\varepsilon$, iteration yields the bootstrap.

\paragraph{Step 3: Local slices for the symmetry action.}  \open\quad
At a solution $\rho$, construct a complement $S_\rho$ to the orbit tangent
\[
  \mathfrak g\cdot\rho
  = \{-i(a\cdot\omega)\rho-(A\omega)\cdot\nabla_{S^{n-1}}\rho:
  a\in\R^n,\;A\in\mathfrak{so}(n)\}.
\]
This must be done separately on orbit-type strata (radial, zonal, fully
asymmetric solutions have different stabiliser dimensions).

\paragraph{Step 4: Nondegeneracy modulo symmetry.}  \open\quad
Verify $\ker L_\rho=\mathfrak g\cdot\rho$ at each solution, where
$L_\rho=I-ckA_k(\rho,\ldots,\rho,\cdot)$.  For the $S^1$ problem at
the constant solution, this reduces to Bessel integrals
(Theorem~\ref{thm:nondeg-criterion}).

\paragraph{Step 5: Compactness modulo~$G$.}  \open\quad
Prove sequential compactness of the solution set modulo~$G$:
\[
  \rho_j=cT(\rho_j)\implies\exists g_j\in G:\;g_j\cdot\rho_j\to\rho_\infty.
\]
This requires a priori bounds ($\norm\rho\leq R$), a non-collapse lower
bound ($\norm\rho\geq r>0$), and a no-escape condition for translations.
Rotations are compact; translations are the genuine noncompact obstruction.

\subsection{Degenerate fibres}
\label{sec:degenerate-fibres}

The fibre $V_k(\omega)=\{(\omega_1,\ldots,\omega_k)\in(S^{n-1})^k:
\sum\omega_j=\omega\}$ is generically a smooth manifold of dimension
$(n-1)k-(n-1+1)=(n-1)(k-1)-1$... when the differential of
$\Sigma:\omega_1+\cdots+\omega_k$ has full rank.  Degeneracies occur when
all tangent lines $T_{\omega_j}S^{n-1}$ are parallel.

\textbf{Example for $k=5$ on $S^1$:}  The configuration
$(\omega_0,\omega_0,\omega_0,-\omega_0,-\omega_0)$ with output
$\omega=\omega_0$ is degenerate.  Here
$d\Sigma=3(\text{tangent at }\omega_0)-2(\text{tangent at }-\omega_0)$
has rank~$1$ instead of~$2$ (the tangents at $\omega_0$ and $-\omega_0$
are parallel on $S^1$).  This is a codimension-$1$ degeneracy, producing
an integrable singularity $\sim|t|^{-1/2}$ in the fibre density.

In higher dimensions, the degenerate locus can have higher codimension
(tangents at $\omega_0$ and $-\omega_0$ are parallel in $\R^n$ only
when $n=1$; for $n\geq 2$, anti-parallel tangent \emph{spaces} are
generically transverse).  This suggests that the smoothing estimate may
be easier for $n\geq 3$ than for $n=2$.

\subsection{Orbit dimension and radial rigidity}

\begin{proposition}[Orbit dimension, $n=3$]
  For $G=\R^3\rtimes SO(3)$ ($\dim G=6$), the orbit dimension of a
  solution~$u$ is $6-\dim\Stab_G(u)$.  Specifically:
  \begin{itemize}
    \item Radial: $d=3$ (only translations).
    \item Zonal (one axis of symmetry): $d=5$.
    \item Fully asymmetric: $d=6$.
  \end{itemize}
\end{proposition}

\begin{conjecture}[Radial rigidity]
  \label{conj:radial}
  For the equation $u=c(|u|^{2(k-1)}u\ast\wh\sigma)$ on $\R^3$, every
  non-trivial solution is rotationally invariant about some point (its
  ``centre of mass'').  If true, all solutions belong to $3$-dimensional
  translation orbits, and the shape count reduces to counting radial profiles.
\end{conjecture}

\begin{warning}
  Conjecture~\ref{conj:radial} is currently \emph{unproved}.  The
  preceding sections on the $n=3$ problem treated it as established; it
  should be regarded as a useful heuristic until proved.
\end{warning}

% =======================================================================
\section{Supplementary: Gegenbauer Expansion and Galerkin Models}
\label{sec:galerkin}
% =======================================================================

\subsection{Gegenbauer expansion}
\heuristic

For completeness, we record the formal harmonic decomposition.
Let $\rho(\omega)=\sum_\ell a_\ell C_\ell^\lambda(\omega\cdot\eta)$ with
$\lambda=(n-2)/2$.  The product of Gegenbauer polynomials linearises as
$C_{\ell_1}^\lambda C_{\ell_2}^\lambda
=\sum_j\gamma_j(\ell_1,\ell_2)C_j^\lambda$.
If $T$ were a spherical convolution, the Funk--Hecke theorem would give
$C_\ell^\lambda\ast_S\sigma=\lambda_\ell(\sigma)C_\ell^\lambda$.
Substituting into $\rho=cT(\rho)$ and equating coefficients would give an
algebraic system for $\{a_\ell\}$.

\textbf{However}, as noted in Warning~\ref{warn:not-convolution}, the
incidence-fibre operator is \emph{not} spherical convolution.  The
Gegenbauer expansion gives a useful \emph{model} but does not rigorously
reduce to finitely many coefficients without separate justification.

\subsection{Galerkin approximation}

Let $H_L=\operatorname{span}\{Y_{\ell m}:0\leq\ell\leq L\}$ and $P_L$ the
$L^2(S^{n-1})$ projection.  The projected equation $P_LF(\rho_L)=0$ is a
polynomial system of degree~$k$ in $N(L)=(L+1)^2$ unknowns (for $n=3$).

\begin{proposition}[Galerkin shape count]
  After imposing algebraic slice conditions, the number of Galerkin shapes
  at truncation~$L$ is bounded by the BKK mixed volume of the sliced
  Newton polytopes, in particular by $k^{N(L)}$ (B\'ezout).
\end{proposition}

\begin{warning}[Limit caveat]
  This is a statement about the fixed-$L$ model.  It does not imply
  finiteness of true shapes as $L\to\infty$ without uniform a priori bounds,
  no-escape compactness, and stability of roots under high-mode
  perturbation.
\end{warning}

\begin{remark}[How Galerkin helps]
  A Galerkin computation becomes rigorous if paired with a high-mode bound
  $\norm{(I-P_L)T(\rho)}_X\leq\varepsilon_L\Phi(\norm\rho_X)$ with
  $\varepsilon_L\to 0$, plus an inverse bound on the sliced Jacobian at each
  computed root.  A Newton--Kantorovich argument would then validate nearby
  true solutions and certify completeness.
\end{remark}

% =======================================================================
\section{Supplementary: Validated Search Strategy}
\label{sec:validated-search}
% =======================================================================

Assuming the transversality route is completed (modulo~$G$ the solution set
is finite), the next goal is to \emph{find every shape} by a finite,
checkable search.

\subsection{Constants needed}

\begin{enumerate}[leftmargin=2em]
  \item \textbf{A priori radius:} $\norm\rho_X\leq R(c,k)$.
  \item \textbf{Non-collapse:} $\norm\rho_X\geq r(c,k)>0$.
  \item \textbf{Regularity/mesh modulus:} $\norm\rho_{C^m}\leq B_m$.
  \item \textbf{Validated quadrature:} error $q_h(B_m)\to 0$ for the
    fibre integral on a mesh of size~$h$.
  \item \textbf{Slice conditioning:}
    $\norm{L_\rho\eta}_X\geq\mu(\rho)\norm\eta_X$ on $S_\rho$, with
    $\mu_*:=\inf_\rho\mu(\rho)>0$.
  \item \textbf{Lipschitz:}
    $\norm{DF(\rho)-DF(\tilde\rho)}\leq M\norm{\rho-\tilde\rho}_X$.
\end{enumerate}

\subsection{Separation and packing}

From the conditioning bound, the quantitative isolation scale is
$\delta_*\simeq\mu_*/M$: two distinct shapes cannot lie within
$\delta_*$ of each other in the slice.  A crude packing bound is
\[
  \#\{\text{shapes}\}
  \leq
  \left(1+\frac{2R_h}{\delta_*-2\eta_h-2q_h}\right)^{N_{\mathrm{mesh}}(h)},
\]
provided $\delta_*>2\eta_h+2q_h$.

\subsection{The search pipeline}

\begin{enumerate}[leftmargin=2em]
  \item Choose a quasi-uniform mesh on $S^{n-1}$ (or directly on $S^1$
    for the circle problem) and fix gauge conditions.
  \item Represent candidates by interval boxes for nodal values plus global
    regularity bounds.
  \item For each box, enclose the residual $F(\rho)$ at mesh points using
    validated quadrature.
  \item Discard boxes with nonzero residual enclosure.
  \item Apply Krawczyk/interval Newton to surviving boxes.
  \item Validate uniqueness and lift to the full space.
  \item Certify completeness: every box discarded or assigned.
\end{enumerate}

\subsection{Main difficulties}

\begin{enumerate}[leftmargin=2em]
  \item \textbf{Translation gauge:} $e^{-ix\cdot\omega}\rho(\omega)$ creates
    oscillation as $|x|\to\infty$.  A canonical centre must be fixed.
  \item \textbf{Uniform conditioning:} $\mu_*>0$ may be the hardest
    constant to obtain without knowing all solutions.
  \item \textbf{Mesh consistency:} nodal errors must not hide solutions
    between mesh points.
  \item \textbf{Dimension:} the effective search dimension grows with
    mesh size.  Symmetry decomposition and adaptive refinement are essential.
\end{enumerate}

\subsection{Specialisation to $S^1$}

For the circle problem, the validated search simplifies considerably:
\begin{itemize}
  \item The mesh is a grid on $[0,2\pi)$ with $N$ equally spaced points.
  \item The gauge fixes phase, rotation of $S^1$, and one translation.
  \item The fibre integral reduces to a trigonometric integral.
  \item The Jacobian involves the linearised operator~$L$, which decomposes
    into Fourier modes for radial $u_0$.
\end{itemize}

The direct certified programme for the sharp constant is:
\begin{enumerate}[leftmargin=2em]
  \item Prove a priori bounds:
    $\norm{f}_{L^2}=1$, $\norm{f}_{C^m}\leq B_m$,
    $0<\Lambda_{\min}\leq\Lambda\leq C_{\mathrm{ST}}^6$.
  \item Use Lemma~\ref{lem:positive-reduction} to restrict the extremiser
    search to boxes containing a nonnegative representative.
  \item Fix gauges (phase, rotation, translation); in the nonnegative
    slice, Lemma~\ref{lem:positive-gauge} fixes the translation parameter.
  \item Mesh $S^1$; represent $f$ by interval boxes plus derivative bounds.
  \item For each $(f,\Lambda)$ box, enclose the residual
    $E^*(|Ef|^4Ef)(\omega_i)-\Lambda f(\omega_i)$.
  \item Discard or validate via interval Newton.
  \item Sort validated $\Lambda_j$ intervals; certify the ordering.
\end{enumerate}

% =======================================================================
\section{Summary: The Logical Dependency Graph}
\label{sec:dependency}
% =======================================================================

We summarise the dependencies and status of the full programme.

\bigskip
\begin{center}
\fbox{\parbox{0.9\textwidth}{
\textbf{Proved:}
\begin{itemize}[leftmargin=1.5em, topsep=2pt, itemsep=1pt]
  \item Spectral support forcing ($\wh u\subseteq S^1$)
  \item 1D case complete
  \item Compactness of linearised operator $K_{u_0}$ on $L^2(S^1)$
  \item Fredholm index zero for $L=K-\Lambda I$
  \item Eigenvalue decomposition at the constant solution
  \item Nondegeneracy $\Leftrightarrow$ Bessel integrals~\eqref{eq:nondeg-cond}
  \item \textbf{Nondegeneracy at the constant: complete proof}
    (Theorem~\ref{thm:nondeg-proved}): Arb certificate for
    $2\leq n\leq12$, explicit Landau/Krasikov bound for $n\geq13$
  \item Spectral gap $\nu_0\geq 0.4\Lambda_0$ certified
  \item Regularity bootstrap for critical points
  \item Inverse Strichartz inequality for $S^1$ extension
  \item Profiles of critical sequences are critical points
  \item Existence of a maximiser (Theorem~\ref{thm:maximiser})
  \item Positivity reduction for extremisers
    (Lemma~\ref{lem:positive-reduction})
  \item Finiteness of nonnegative critical values in compact positive
    windows (Theorem~\ref{thm:positive-values-finite})
  \item NK isolation radius $r_0\geq 0.13$ (floating; Arb pending)
\end{itemize}

\review{The last item should not sit under ``Proved'' in a final version:
it is explicitly floating-point and Arb-pending.  A better status split is
``proved'', ``partially certified'', and ``numerical target''.  Likewise,
the sharp-constant route currently needs a quantitative BTZ stability
certificate and a rigorous high-mode interaction bound before the
conditional conclusion becomes a proof.}

\response{\resolved\ Agreed.  The status box will be revised to use
three categories: \textbf{Proved} (complete proofs), \textbf{Partially
certified} (correct arguments with floating-point constants pending Arb),
and \textbf{Numerical target} (computations not yet attempted).  The NK
radius moves to ``Partially certified''.  The sharp-constant conditional
conclusion will list all four pending certificates explicitly:
(1)~Bessel gap, (2)~NK radius, (3)~high-mode cross bound,
(4)~quantitative BTZ stability modulus.}

\textbf{Next concrete steps:}
\begin{itemize}[leftmargin=1.5em, topsep=2pt, itemsep=1pt]
  \item Certify $M_0$ and $r_0$ by Arb ball arithmetic
    (sixfold Bessel integrals)
  \item Connect the band-limited uniqueness of Barker--Thiele--ZK
    (modes $\leq 120$) with the bootstrap tail to bypass the far-region
    branch-and-bound search
  \item Prove the constant is the unique maximiser modulo $G$
    (this alone gives $\mathcal R$)
\end{itemize}

\textbf{Open (harder):}
\begin{itemize}[leftmargin=1.5em, topsep=2pt, itemsep=1pt]
  \item Rule out multi-profile splitting for non-maximising
    critical sequences (Theorem~\ref{thm:compactness-mod-G})
  \item Enumeration / uniqueness of non-constant shapes
\end{itemize}

\textbf{Conditional conclusion:}
\begin{itemize}[leftmargin=1.5em, topsep=2pt, itemsep=1pt]
  \item Compactness + nondegeneracy $\Longrightarrow$ finitely many
    critical values $\Lambda_j$
  \item $\mathcal R=\max_j\Lambda_j^{1/6}$
\end{itemize}
}}
\end{center}

\bigskip
\noindent\rule{\linewidth}{0.4pt}
\begin{center}
  \small\itshape
  Notes compiled from working discussions.  Statements marked \proven{}
  carry complete proofs; statements marked \conditional{} depend on
  explicitly stated hypotheses; statements marked \open{} require new
  analytic input.
\end{center}

% =======================================================================
\section{Working notes: failed and partial approaches}
\label{sec:working-notes}
% =======================================================================

This section records approaches that were explored but did not succeed,
together with the specific obstructions encountered.  The material is
kept for reference and to avoid repeating dead ends.

\subsection{Approaches to bounding the spectral gap $\gamma_0$}

The spectral gap $\gamma_0:=\inf\{\Lambda-\Lambda_0:\Lambda\in
\mathcal V_+,\,\Lambda>\Lambda_0\}$ measures the distance from the
constant's eigenvalue to the next nonneg critical value above it.
The finiteness theorem (Theorem~\ref{thm:positive-values-finite})
gives $\gamma_0>0$ non-constructively.  We attempted several routes
to an effective lower bound; none succeeded cleanly.

\paragraph{Cauchy--Schwarz envelope.}
For nonneg $f$ with $\norm{f}_2=1$, write
$|Ef(r,\phi)|\leq 2\pi a_0|J_0(r)|+2\pi\sigma_0 S(r)$
where $S(r):=\sqrt{\sum_{n\geq 2}J_n(r)^2}$.
By the Neumann addition theorem,
$S(r)^2=(1-J_0(r)^2-2J_1(r)^2)/2$, which is $O(1)$ for large~$r$
(no decay).  Consequently $\int[a_0|J_0|+\sigma_0 S]^6 r\,dr$ diverges
for $\sigma_0>0$, and the envelope bound cannot certify
$J(f)<\Lambda_0$.  \textbf{Obstruction:} loss of oscillatory
cancellation between Bessel modes.

\paragraph{Leading-order sextic expansion.}
Grouping the sextuples by number of nonzero modes:
$J(f)=\Lambda_0(2\pi a_0^2)^3+C_2 a_0^4\sigma_0^2+C_4 a_0^2\sigma_0^4
+C_6\sigma_0^6$,
with $C_2=15\,\mathcal J_2/\mathcal J_0\approx 1.65$ from the
$(4\text{ zero},2\text{ nonzero})$ sextuples.  At leading order
this gives $\gamma_0\gtrsim 12$.  However, the coefficients $C_4,C_6$
involve Bessel integrals with $4$--$6$ nonzero indices, whose bounds
require the same type of computation as $c_{\mathrm{stab}}$ itself.
The leading-order estimate is therefore not self-contained.

\paragraph{Truncation of the EL residual.}
Attempting to certify $T(f)-\Lambda f\neq 0$ mode-by-mode with a
Fourier truncation at level~$N_0$: the Lipschitz constant of~$T$ is
$\sim 10\,C_{\mathrm{ST}}^6\approx 7000$, while the mode-$0$ signal
$(T(f))_0\sim\Lambda_0 a_0\approx 210$.  The truncation error exceeds
the signal for \emph{every}~$N_0$, because $\sigma_{N_0}\geq
c/(2\pi)^{1/2}$ for nonneg~$f$ far from~$f_0$.  \textbf{Obstruction:}
the quintic nonlinearity has too large a Lipschitz constant relative
to the eigenvalue spacing.

\paragraph{BTZ SOS eigenvalues.}
The BTZ proof reduces to positive semi-definiteness of matrices $Q_D$
whose entries are sixfold Bessel integrals.  The minimum eigenvalue
$\lambda_{\min}(Q_0)\sim 0.15\,N^{-1.74}$ (Figure~2 of BTZ-II)
certifies $\Lambda_0-J(f)\geq\lambda_{\min}\sum|v_{\mathbf m}|^2$
where $v_{\mathbf m}=\wh f_{m_1}\wh f_{m_2}\wh f_{m_3}$ are
triple products.  Converting to a quadratic stability
$c_{\mathrm{stab}}\,d_G^2$ introduces a factor $\wh f_0^4\sim
(2\pi)^{-2}$, giving $c_{\mathrm{stab}}\sim 10^{-6}$---far
below the needed threshold.  \textbf{Obstruction:} the SOS basis
(cubic products) is not adapted to quadratic distance from~$f_0$.

\subsection{The closure argument and its cost}

The approach that \emph{does} work, at least in principle, is the
$c_{\mathrm{stab}}$ closure from
Section~\ref{sec:btk-connection}: certify
\[
  c_{\mathrm{stab}}(N)
  :=
  \inf_{\substack{g\text{ real, nonneg, }\norm{g}_2=1\\
  \wh g(n)=0\text{ for }|n|>N,\;g\neq f_0\bmod G}}
  \frac{\Lambda_0-J(g)}{d_G(g,f_0)^2}
  \;>\;
  \mathrm{threshold}(N),
\]
where $\mathrm{threshold}(N):=(C_Q-3)\Lambda_0\sigma_N^2/(r_0-\sigma_N)^2$
with $\sigma_N=A_3/(\sqrt{\Lambda_-}\,N^{5/2})$ from
Corollary~\ref{cor:tail-bound}.

Both sides depend on~$N$.  The threshold decays as $N^{-5}$
(from $\sigma_N^2$).  Floating-point scans give
$c_{\mathrm{stab}}(2)\geq 128$ and $c_{\mathrm{stab}}(4)\geq 39$,
suggesting a decay $c_{\mathrm{stab}}(N)\sim C_0 N^{-\alpha}$ with
$\alpha\approx 1.7$ (fitted from two points; consistent with the BTZ
eigenvalue decay $\sim N^{-1.74}$).

Since $\alpha<5$, the threshold eventually drops below
$c_{\mathrm{stab}}$ for $N$ large enough.  The crossing
point~$N_{\mathrm{cross}}$ and the associated BTZ computational cost
$\sim N_{\mathrm{cross}}^5$ depend on the actual value of~$\alpha$:

\begin{center}
\begin{tabular}{ccc}
\hline
$\alpha$ & $N_{\mathrm{cross}}$ & BTZ cost ($\sim N^5$) \\
\hline
$0$ (no decay) & $\sim 140$ & $\sim 5\times 10^{10}$ \\
$1$ & $\sim 330$ & $\sim 4\times 10^{12}$ \\
$1.74$ (empirical) & $\sim 900$ & $\sim 6\times 10^{14}$ \\
$2$ & $\sim 1400$ & $\sim 6\times 10^{15}$ \\
$3$ & $\sim 27000$ & $\sim 10^{22}$ \\
\hline
\end{tabular}
\end{center}

For comparison, the published BTZ computation at $N=120$ used $94$
compute nodes for $16$~hours with $\sim 1.6\times 10^8$ Bessel
integrals.

\paragraph{Next steps.}  The decay rate~$\alpha$ can be pinned down by
computing $c_{\mathrm{stab}}(N)$ at $N=6,8,10,\ldots$ via the global
stability scan (\texttt{numerics/global\_stability\_v2.py}).  Once
$\alpha$ is known to one decimal place, the required~$N$ is determined,
and the computational cost can be quoted to a numerics team for a
BTZ-type Arb computation at level~$N_{\mathrm{cross}}$.


\end{document}
